\documentclass[onecolumn]{IEEEtran}

% -------------------------------------------------
% Paquetes
% -------------------------------------------------

\usepackage{amsmath}
\usepackage{amssymb}
\usepackage{algorithm}
\usepackage{algorithmic}
\usepackage{array}
\usepackage{booktabs}
\usepackage{cite}
\usepackage{url}
\usepackage[T1]{fontenc}
\usepackage{graphicx}
\usepackage{epstopdf}

% -------------------------------------------------
% Documento
% -------------------------------------------------

\begin{document}

% -------------------------------------------------
% Título
% -------------------------------------------------

\title{
RICAN: Un Marco de Representación Binaria de Tasa Fija
para Alfabetos Finitos Usando Convergentes de Fracciones Continuas
}

\author{
\IEEEauthorblockN{
J. Jes\'us Mart\'inez Palomo
}
\IEEEauthorblockA{
Investigador Independiente\\
Correo: enviarcorreoajesusdelaunam@gmail.com
}
}

\maketitle

% -------------------------------------------------
% Resumen
% -------------------------------------------------

\begin{abstract}

Este artículo presenta RICAN (Representational Isomorphic Coding for
Arbitrary Numeration), un marco matemático para construir
representaciones binarias reversibles de tasa fija asintóticamente
ajustadas para secuencias sobre alfabetos finitos arbitrarios.

La contribución central es un método constructivo para seleccionar
tamaños de bloque mediante aproximación por fracciones continuas de
$\log_2(B)$. El método identifica longitudes de bloque derivadas de
convergentes que generan representaciones deterministas de tasa fija
con brechas de redondeo absolutas arbitrariamente pequeñas a lo largo
de subsecuencias adecuadas. En consecuencia, el costo de
representación normalizado por símbolo converge a cero, donde $B$
denota la cardinalidad del alfabeto.

RICAN no realiza compresión de fuente; minimiza el costo determinista
de representación binaria de tasa fija.

Para un bloque de longitud $K$ sobre un alfabeto de tamaño $B$, el
número mínimo de bits requerido por cualquier representación binaria
inyectiva de longitud fija está acotado por

\[
n(K)=\lceil K\log_2(B)\rceil
\]

bits.

El costo de redondeo inevitable es por lo tanto

\[
G_B(K)=
\lceil K\log_2(B)\rceil-K\log_2(B).
\]

RICAN demuestra que los denominadores de los convergentes de
fracciones continuas de $\log_2(B)$ proporcionan una subsecuencia
infinita de tamaños de bloque cuya brecha de redondeo absoluta tiende
a cero, y por lo tanto cuyo costo de representación normalizado por
símbolo converge a cero.

El marco se aplica uniformemente a la familia infinita de alfabetos
finitos

\[
\mathcal{A}=\{\Sigma_B:B\geq2\},
\]

donde

\[
\Sigma_B=\{0,1,\ldots,B-1\}.
\]

Los alfabetos de potencia de dos logran representaciones binarias
exactas con costo cero. Para alfabetos no potencia de dos, los
convergentes de fracciones continuas proporcionan tamaños de bloque
deterministas asintóticamente ajustados.

Para el alfabeto decimal ($B=10$), la implementación práctica de
referencia utiliza el tamaño de bloque derivado del convergente

\[
K=643,
\]

que requiere

\[
n(K)=2136
\]

bits y logra un costo de representación de aproximadamente

\[
3.65\times10^{-7}
\]

bits por símbolo.

Un denominador convergente superior mayor

\[
K=783032
\]

corresponde al convergente

\[
\frac{2601184}{783032}
\]

en la expansión en fracciones continuas de $\log_2(10)$, satisfaciendo

\[
K\log_2(10)=2601183.99999948\ldots,
\]

requiriendo

\[
n(K)=\lceil K\log_2(10)\rceil=2601184
\]

bits, representados exactamente por

\[
325148
\]

bytes.

La brecha de redondeo total de este bloque completo es aproximadamente
$5.20\times10^{-7}$ bits, mientras que el costo normalizado por
símbolo es

\[
\epsilon_{10}(783032)
=
\frac{5.20\times10^{-7}}{783032}
\approx
6.64\times10^{-13}
\]

bits por símbolo.

En este contexto, la proximidad de representación se refiere
exclusivamente al costo de realización binaria y no implica
compresión de fuente.

A diferencia de los métodos estadísticos de codificación de fuente,
RICAN no explota la redundancia de la fuente ni las distribuciones de
probabilidad. La longitud codificada depende únicamente de la
cardinalidad del alfabeto y la longitud de entrada.

La contribución principal es un procedimiento constructivo para
seleccionar representaciones de bloques de tasa fija con brecha de
redondeo absoluta asintóticamente evanescente mediante aproximación
por fracciones continuas.

Se proporciona una implementación de referencia que utiliza el formato
de representación autónomo RJ03.

\end{abstract}

% -------------------------------------------------
% Términos de Indexación
% -------------------------------------------------

\begin{IEEEkeywords}

Representación de tasa fija,
alfabetos finitos,
fracciones continuas,
aproximación diofántica,
representación enumerativa,
codificación biyectiva.

\end{IEEEkeywords}

% =================================================
% I. INTRODUCCIÓN
% =================================================

\section{Introducción}

Los sistemas de información sin pérdida requieren mapeos reversibles
entre secuencias de símbolos y descripciones binarias.

Dos objetivos diferentes se asocian frecuentemente con tales mapeos:
compresión de fuente y representación binaria.

Los métodos de compresión intentan reducir la longitud esperada de
descripción explotando propiedades estadísticas de la fuente. En
contraste, los métodos de representación estudian el mapeo
determinista entre un alfabeto finito y un espacio de almacenamiento
binario.

RICAN pertenece a la segunda categoría.

El objetivo de RICAN no es comprimir el contenido de información, sino
minimizar el costo de redondeo entero inevitable que aparece cuando
espacios de alfabetos finitos se representan utilizando hardware
binario.

Para un alfabeto que contiene $B$ símbolos, la cantidad de información
asociada con un símbolo seleccionado uniformemente es

\[
\log_2(B)
\]

bits.

Cuando

\[
B=2^m ,
\]

la representación es exacta porque

\[
\log_2(B)=m .
\]

Sin embargo, para cardinalidades de alfabeto que no son potencias de
dos, $\log_2(B)$ es generalmente irracional, impidiendo una
representación binaria exacta de tasa fija de un símbolo.

Una representación por bloques resuelve este problema agrupando
símbolos.

Para una longitud de bloque $K$, el número de secuencias posibles es

\[
|\Sigma_B^K|=B^K .
\]

Una representación binaria que contiene $n$ bits proporciona

\[
2^n
\]

estados posibles.

Por lo tanto, cualquier representación reversible de tasa fija debe
satisfacer

\[
2^n\geq B^K .
\]

Tomando logaritmos se obtiene:

\[
n\geq K\log_2(B).
\]

Debido a que el número de bits debe ser entero, la longitud mínima
requerida del bloque es:

\[
n(K)=\lceil K\log_2(B)\rceil .
\]

El costo de redondeo resultante es:

\[
G_B(K)=
\lceil K\log_2(B)\rceil
-
K\log_2(B).
\]

El problema central abordado por RICAN es la selección de tamaños de
bloque $K$ que minimizan esta brecha determinista.

La contribución de este trabajo es un método constructivo basado en
fracciones continuas.

Para valores irracionales de $\log_2(B)$, los denominadores de los
convergentes de fracciones continuas proporcionan longitudes de
bloque donde el producto

\[
K\log_2(B)
\]

se aproxima a un entero con precisión creciente.

El marco está definido sobre la familia infinita de alfabetos finitos:

\[
\mathcal{A}
=
\{\Sigma_B:B\geq2\}.
\]

Por lo tanto, alfabetos binarios, alfabetos de ADN, dígitos decimales,
símbolos hexadecimales, alfabetos de bytes y alfabetos finitos
arbitrarios son tratados como instancias particulares del mismo modelo
matemático.

Las principales contribuciones de este artículo son:

\begin{enumerate}

\item Un marco matemático general de representación binaria de tasa fija
para alfabetos finitos.

\item Un procedimiento constructivo de selección de tamaños de bloque
basado en convergentes de fracciones continuas de $\log_2(B)$.

\item Un análisis teórico que prueba la existencia de subsecuencias
convergentes con brecha de representación absoluta evanescente.

\item Un formato de representación binaria reversible autónomo (RJ03)
que almacena toda la información requerida para la reconstrucción
determinista.

\item Una validación experimental sobre secuencias decimales y
múltiples alfabetos finitos.

\end{enumerate}

El resto de este artículo está organizado de la siguiente manera.

La Sección II revisa enfoques relacionados.

La Sección III introduce los fundamentos matemáticos.

La Sección IV presenta los principales resultados teóricos.

La Sección V describe el algoritmo de selección de bloques.

La Sección VI presenta el marco de representación RICAN y el formato
RJ03.

La Sección VII presenta la evaluación experimental.

La Sección VIII discute implicaciones y limitaciones.

La Sección IX presenta la disponibilidad del código.

La Sección X concluye el artículo.

% =================================================
% NOTACIÓN
% =================================================

\section{Notación}

La siguiente notación se utiliza a lo largo de este artículo.

\begin{table}[!t]
\caption{Resumen de Notación}
\label{tab:notation}
\centering
\begin{tabular}{ll}
\toprule
Símbolo & Significado\\
\midrule
$B$ & Cardinalidad del alfabeto\\
$\Sigma_B$ & Alfabeto de tamaño $B$\\
$K$ & Longitud de bloque en símbolos\\
$n(K)$ & Longitud de representación binaria de un bloque\\
$G_B(K)$ & Brecha de representación absoluta\\
$\delta_B(K)$ & Costo de representación normalizado\\
$\epsilon_B(K)$ & Costo de representación por símbolo\\
$p_n/q_n$ & Convergente de fracción continua\\
$L$ & Longitud de la secuencia de entrada\\
$N(L)$ & Longitud total de representación\\
\bottomrule
\end{tabular}
\end{table}

La notación enfatiza la distinción entre tamaño de fuente y tamaño de
representación.

La cardinalidad del alfabeto $B$ determina la tasa ideal
$\log_2(B)$, mientras que la longitud de bloque $K$ determina el
comportamiento de redondeo entero.

% =================================================
% II. TRABAJO RELACIONADO
% =================================================

\section{Trabajo Relacionado}

El problema de representar información de alfabetos finitos utilizando
hardware binario ha sido estudiado desde varias perspectivas,
incluyendo codificación estadística, representación enumerativa y
técnicas de mapeo de longitud fija.

RICAN difiere de estos enfoques porque su objetivo no es la reducción
de la longitud esperada de descripción de la fuente, sino la
minimización del costo determinista de representación binaria.

\subsection{Codificación Estadística de Fuente}

Los métodos clásicos de codificación de fuente optimizan la longitud
de representación cuando las propiedades estadísticas de la fuente
están disponibles.

La codificación de Huffman asigna palabras de código prefijo de
longitud variable según las probabilidades de los símbolos y
proporciona optimalidad entre los códigos prefijo binarios para una
distribución conocida \cite{huffman1952}.

La codificación aritmética representa secuencias como subintervalos del
intervalo unitario y evita la limitación de longitud entera de la
codificación prefijo convencional \cite{rissanen1979}.

Los sistemas numéricos asimétricos (ANS) logran un rendimiento de
codificación entrópica comparable a la codificación aritmética
permitiendo implementaciones eficientes con máquinas de estados finitos
\cite{duda2013}.

Estas técnicas explotan la redundancia estadística.

Su objetivo es por lo tanto diferente de RICAN, que construye
representaciones de tasa fija cuyo tamaño depende únicamente de la
cardinalidad del alfabeto y la longitud de la secuencia.

\subsection{Métodos Basados en Diccionario}

Los métodos basados en diccionario reducen la longitud de descripción
de la fuente identificando estructuras repetidas en las secuencias.

La familia Lempel-Ziv introdujo algoritmos universales que representan
subcadenas repetidas mediante referencias a patrones observados
previamente \cite{ziv1977}.

Los sistemas de compresión prácticos modernos frecuentemente combinan
métodos de diccionario con codificación entrópica para mejorar el
rendimiento en datos estructurados.

RICAN no realiza detección de patrones, estimación de probabilidades o
eliminación de redundancia.

En cambio, aborda el problema independiente de representar cada bloque
posible de un alfabeto finito mediante un mapeo binario determinista
con costo mínimo de redondeo entero.

\subsection{Representación Enumerativa}

Los métodos enumerativos proporcionan procedimientos sistemáticos de
clasificación y desclasificación para espacios combinatorios finitos.

Cover introdujo la codificación enumerativa de fuente como un método
para mapear conjuntos de secuencias a índices enteros
\cite{cover1973}.

Desarrollos posteriores extendieron estas ideas a secuencias
restringidas y representaciones combinatorias eficientes
\cite{immink2004}.

RICAN está relacionado con la representación enumerativa porque ambos
enfoques utilizan la cardinalidad de espacios finitos.

Sin embargo, el objetivo principal difiere.

La codificación enumerativa se centra en clasificar clases específicas
de secuencias, frecuentemente bajo restricciones.

RICAN se centra en la selección de longitudes de bloque $K$ tales que
la longitud de representación binaria

\[
\lceil K\log_2(B)\rceil
\]

tenga un error de redondeo determinista mínimo.

\subsection{Representación Binaria de Longitud Fija}

Un alfabeto finito con cardinalidad $B$ requiere al menos

\[
\lceil\log_2(B)\rceil
\]

bits cuando los símbolos se representan independientemente.

Este enfoque es óptimo sólo cuando $B$ es una potencia de dos.

Para alfabetos generales, la representación por bloques reduce el
efecto del redondeo considerando el espacio de bloques completo:

\[
|\Sigma_B^K|=B^K .
\]

La longitud binaria mínima es por lo tanto:

\[
n(K)=\lceil K\log_2(B)\rceil .
\]

RICAN estudia la selección matemática de $K$ en lugar de la
construcción de un nuevo modelo probabilístico de codificación.

\subsection{Fracciones Continuas y Aproximación Diofántica}

Las fracciones continuas proporcionan aproximaciones racionales
altamente precisas a números irracionales.

Para un número irracional $x$, la expansión en fracciones continuas es

\[
x=[a_0;a_1,a_2,\ldots].
\]

Los convergentes son fracciones

\[
\frac{p_n}{q_n}
\]

con la propiedad de aproximación clásica:

\[
\left|
x-\frac{p_n}{q_n}
\right|
<
\frac{1}{q_nq_{n+1}} .
\]

Multiplicando por $q_n$ se obtiene:

\[
|q_nx-p_n|
<
\frac{1}{q_{n+1}} .
\]

Por lo tanto, los denominadores $q_n$ identifican múltiplos enteros de
$x$ que se aproximan a valores enteros con precisión creciente.

RICAN aplica esta propiedad clásica de teoría de números a la
selección de tamaños de bloque de representación.

Para

\[
x=\log_2(B),
\]

los denominadores seleccionados proporcionan longitudes de bloque
donde:

\[
K\log_2(B)
\]

está cerca de un límite entero.

\subsection{Posición de RICAN}

RICAN debe interpretarse como un marco matemático de representación.

Su contribución no es un algoritmo de codificación entrópica ni un
mecanismo de compresión.

El marco proporciona:

\begin{itemize}

\item un modelo general para representación de alfabetos finitos,

\item un método para seleccionar tamaños de bloque de tasa fija,

\item un análisis teórico basado en aproximación por fracciones
continuas,

\item un formato binario reversible autónomo.

\end{itemize}

La distinción conceptual entre RICAN y los métodos de codificación
convencionales se resume en la Tabla~\ref{tab:comparison}.

\begin{table}[!t]
\caption{Comparación Conceptual Entre Enfoques de Representación y
Compresión}
\label{tab:comparison}
\centering
\begin{tabular}{lll}
\toprule
Método & Objetivo & Estadísticas Requeridas\\
\midrule
Codificación Huffman &
Minimización de longitud esperada &
Sí\\

Codificación Aritmética &
Aproximación de tasa entrópica &
Sí\\

ANS &
Eficiencia de codificación entrópica &
Sí\\

Métodos de Diccionario &
Reducción de redundancia de patrones &
Sí\\

Codificación Enumerativa &
Clasificación de espacio finito &
A veces\\

RICAN &
Minimización de brecha de representación de tasa fija &
No\\

\bottomrule
\end{tabular}
\end{table}

% =================================================
% III. PRELIMINARES
% =================================================

\section{Preliminares}

Esta sección introduce las definiciones matemáticas requeridas para el
marco RICAN.

\subsection{Familia de Alfabetos Finitos}

RICAN considera la familia de alfabetos finitos:

\[
\mathcal{A}
=
\{\Sigma_B:B\geq2\},
\]

donde:

\[
\Sigma_B=
\{0,1,\ldots,B-1\}.
\]

El valor $B$ determina completamente el problema de representación.

Ejemplos representativos incluyen:

\begin{itemize}

\item $B=2$: alfabeto binario.

\item $B=4$: alfabeto de ADN.

\item $B=10$: dígitos decimales.

\item $B=16$: alfabeto hexadecimal.

\item $B=256$: alfabeto de bytes.

\end{itemize}

\subsection{Racionalidad de $\log_2(B)$}

\begin{lemma}
\label{lem:rationality}

Para un entero $B\geq2$,

\[
\log_2(B)\in\mathbb{Q}
\]

si y sólo si

\[
B=2^m
\]

para algún entero $m\geq1$.

\end{lemma}

\begin{proof}

Si:

\[
B=2^m,
\]

entonces:

\[
\log_2(B)=m,
\]

que es racional.

Recíprocamente, supongamos:

\[
\log_2(B)=\frac{p}{q}
\]

en términos mínimos.

Entonces:

\[
B^q=2^p .
\]

Por factorización única, $B$ puede contener únicamente el factor primo
$2$. Por lo tanto:

\[
B=2^m
\]

para algún entero $m$.

\end{proof}

\subsection{Espacio de Bloques}

Para una longitud de bloque $K$, el conjunto de bloques de entrada
posibles es:

\[
\Sigma_B^K .
\]

Su cardinalidad es:

\[
|\Sigma_B^K|=B^K .
\]

Una palabra binaria de longitud $n$ proporciona:

\[
2^n
\]

estados posibles.

Por lo tanto, una representación inyectiva requiere:

\[
2^n\geq B^K .
\]

Tomando logaritmos:

\[
n\geq K\log_2(B).
\]

Debido a que $n$ debe ser entero:

\[
n(K)=\lceil K\log_2(B)\rceil .
\]

\subsection{Brecha de Representación Absoluta}

\begin{definition}

Para una cardinalidad de alfabeto $B$, la brecha de representación
absoluta es:

\[
G_B(K)=
\lceil K\log_2(B)\rceil
-
K\log_2(B).
\]

Esta cantidad mide la distancia de redondeo entero, en bits, entre
la longitud de representación ideal de valor real y la longitud real
del bloque binario.

\end{definition}

\subsection{Costo de Representación Normalizado}

\begin{definition}

El costo de representación normalizado es:

\[
\delta_B(K)
=
\frac{G_B(K)}
{K\log_2(B)} .
\]

El costo de representación por símbolo es:

\[
\epsilon_B(K)
=
\frac{G_B(K)}{K}.
\]

Estas cantidades describen la cercanía de la representación relativa
al límite ideal de tasa fija.

No representan eficiencia de compresión.

\end{definition}

\subsection{Representación Asintóticamente Ajustada}

\begin{definition}

Una secuencia de tamaños de bloque $K_i$ define una representación
asintóticamente ajustada con respecto al costo de redondeo si:

\[
\lim_{i\rightarrow\infty}
\delta_B(K_i)=0 .
\]

Para RICAN, la secuencia seleccionada se obtiene de denominadores
convergentes adecuados de fracciones continuas de $\log_2(B)$.

\end{definition}

\subsection{Longitud de Representación de Tasa Fija}

Para una secuencia de entrada de longitud $L$, la longitud total de
representación es:

\[
N(L)=
\left\lceil\frac{L}{K}\right\rceil n(K),
\]

donde:

\[
n(K)=\lceil K\log_2(B)\rceil .
\]

Esta expresión asume que el bloque incompleto final se codifica
usando su longitud de bits real y excluye los metadatos constantes
del contenedor.

Para $K$ grande y $L$ grande, la tasa efectiva se aproxima a:

\[
\log_2(B)
\]

bits por símbolo a lo largo de la subsecuencia convergente
seleccionada.

\subsection{Relación Matemática con la Codificación Enumerativa}

RICAN y la codificación enumerativa comparten el uso de clasificación
entera sobre espacios finitos.

En la codificación enumerativa, las secuencias se mapean a índices
enteros mediante conteo combinatorio.

En RICAN, el mapeo se define por conversión de base sobre bloques
completos:

\[
X(S)=\sum_{i=0}^{K-1}s_iB^{K-1-i}.
\]

Ambos enfoques producen mapeos reversibles.

Sin embargo, el criterio de optimización difiere.

La codificación enumerativa típicamente minimiza la longitud esperada
del índice bajo restricciones de probabilidad.

RICAN minimiza el costo de redondeo determinista a nivel de bloque
mediante selección de bloques por fracciones continuas.

Por lo tanto, RICAN puede interpretarse como un caso especial de
representación enumerativa donde:

\begin{itemize}

\item el espacio enumerado es el producto cartesiano completo
$\Sigma_B^K$,

\item el objetivo es la optimización de la longitud de bloque,

\item la selección de $K$ se realiza mediante aproximación
diofántica.

\end{itemize}

Esta conexión aclara la posición de RICAN dentro de las técnicas de
representación existentes.

% =================================================
% IV. PRINCIPALES RESULTADOS TEÓRICOS
% =================================================

\section{Principales Resultados Teóricos}

Esta sección establece el fundamento matemático de RICAN.

La observación central es que el costo de representación de tasa fija
depende de la distancia entre:

\[
K\log_2(B)
\]

y el valor entero más cercano.

Para una cardinalidad de alfabeto $B$, definimos:

\[
x_B=\log_2(B).
\]

La brecha de representación es:

\[
G_B(K)
=
\lceil Kx_B\rceil-Kx_B .
\]

El objetivo de la construcción RICAN es seleccionar longitudes de
bloque $K$ para las cuales esta cantidad se vuelve arbitrariamente
pequeña.

\subsection{Distancia a Valores Enteros}

Para un número real $y$, definimos su distancia al entero más cercano
como:

\[
\operatorname{dist}(y,\mathbb{Z})
=
\min_{m\in\mathbb{Z}}|y-m|.
\]

Para:

\[
x_B=\log_2(B),
\]

definimos:

\[
d_B(K)
=
\operatorname{dist}(Kx_B,\mathbb{Z}).
\]

Esta cantidad mide qué tan cerca está la longitud de representación
ideal de un número entero de bits.

Cuando el entero más cercano corresponde a la operación de techo:

\[
\lceil Kx_B\rceil,
\]

la brecha de representación satisface:

\[
G_B(K)=d_B(K).
\]

\subsection{Aproximación por Fracciones Continuas}

Sea $x$ irracional y sea:

\[
\frac{p_n}{q_n}
\]

un convergente de su expansión en fracciones continuas.

Una propiedad clásica de las fracciones continuas es:

\[
\left|
x-\frac{p_n}{q_n}
\right|
<
\frac{1}{q_nq_{n+1}} .
\]

Multiplicando por $q_n$ se obtiene:

\[
|q_nx-p_n|
<
\frac{1}{q_{n+1}} .
\]

Debido a que $p_n$ es un entero:

\[
\operatorname{dist}(q_nx,\mathbb{Z})
\leq
|q_nx-p_n|.
\]

Por lo tanto:

\[
\operatorname{dist}(q_nx,\mathbb{Z})
<
\frac{1}{q_{n+1}} .
\]

\begin{lemma}
\label{lem:convergent}

Sea $x$ irracional y sea

\[
\frac{p_n}{q_n}
\]

un convergente de fracción continua de $x$.

Entonces:

\[
\operatorname{dist}(q_nx,\mathbb{Z})
<
\frac{1}{q_{n+1}} .
\]

\end{lemma}

\begin{proof}

El teorema de aproximación clásico para fracciones continuas da:

\[
|q_nx-p_n|
<
\frac{1}{q_{n+1}} .
\]

Como $p_n$ pertenece a $\mathbb{Z}$,

\[
\operatorname{dist}(q_nx,\mathbb{Z})
\leq
|q_nx-p_n|.
\]

Combinando ambas expresiones se prueba el resultado.

\end{proof}

\subsection{Selección de Bloques Basada en Convergentes}

El lema anterior proporciona directamente una regla de construcción
para RICAN.

Para:

\[
x_B=\log_2(B),
\]

elegimos tamaños de bloque:

\[
K=q_n,
\]

donde $q_n$ son denominadores de convergentes de fracciones continuas.

Entonces:

\[
Kx_B
\]

se aproxima a un valor entero con un error acotado por:

\[
\frac{1}{q_{n+1}} .
\]

\begin{theorem}
\label{thm:convergent}

Sea:

\[
x_B=\log_2(B)
\]

irracional.

Sea:

\[
\frac{p_n}{q_n}
\]

los convergentes de la expansión en fracciones continuas de $x_B$.

Entonces las longitudes de bloque:

\[
K_n=q_n
\]

satisfacen:

\[
\operatorname{dist}(K_nx_B,\mathbb{Z})
<
\frac{1}{q_{n+1}} .
\]

En consecuencia, existe una secuencia infinita de tamaños de bloque
RICAN cuya longitud de representación ideal se aproxima a un valor
entero con error absoluto arbitrariamente pequeño.

\end{theorem}

\begin{proof}

Del Lema~\ref{lem:convergent}:

\[
\operatorname{dist}(q_nx_B,\mathbb{Z})
<
\frac{1}{q_{n+1}} .
\]

Los denominadores de los convergentes de fracciones continuas
satisfacen:

\[
q_n\rightarrow\infty .
\]

Por lo tanto:

\[
\frac{1}{q_{n+1}}\rightarrow0 .
\]

Por lo tanto:

\[
\lim_{n\rightarrow\infty}
\operatorname{dist}(q_nx_B,\mathbb{Z})
=0 .
\]

Esto prueba que los denominadores convergentes proporcionan una
secuencia infinita de tamaños de bloque con distancia evanescente
entre las longitudes de representación ideal y entera.

\end{proof}

\subsection{Existencia de Brecha de Representación Absoluta Evanescente}

El teorema anterior se refiere a la distancia al entero más cercano.

Para RICAN, la cantidad relevante es la brecha de techo:

\[
G_B(K)
=
\lceil Kx_B\rceil-Kx_B .
\]

\begin{theorem}
\label{thm:abs_overhead}

Sea:

\[
x_B=\log_2(B)
\]

irracional.

Entonces existe una subsecuencia infinita de denominadores
convergentes de fracciones continuas:

\[
K_i=q_{n_i}
\]

tal que:

\[
G_B(K_i)\rightarrow0 .
\]

Por lo tanto, el costo de representación absoluto de tasa fija de
RICAN puede hacerse arbitrariamente pequeño mediante selección de
bloques basada en convergentes.

\end{theorem}

\begin{proof}

Del Teorema~\ref{thm:convergent}, los denominadores convergentes
satisfacen:

\[
\operatorname{dist}(q_nx_B,\mathbb{Z})
<
\frac{1}{q_{n+1}} .
\]

Los convergentes alternan alrededor del número irracional $x_B$.
Por lo tanto, infinitos convergentes se aproximan a $x_B$ desde
arriba e infinitos desde abajo.

Consideremos la subsecuencia que se aproxima desde arriba:

\[
\frac{p_{n_i}}{q_{n_i}}>x_B .
\]

Para índices suficientemente grandes:

\[
p_{n_i}
=
\lceil q_{n_i}x_B\rceil .
\]

Por lo tanto:

\[
G_B(q_{n_i})
=
p_{n_i}-q_{n_i}x_B .
\]

Usando la cota de fracciones continuas:

\[
G_B(q_{n_i})
<
\frac{1}{q_{n_i+1}} .
\]

Como:

\[
q_{n_i+1}\rightarrow\infty ,
\]

obtenemos:

\[
\lim_{i\rightarrow\infty}
G_B(q_{n_i})
=
0 .
\]

Por lo tanto, la subsecuencia convergente seleccionada proporciona
tamaños de bloque con brecha de representación absoluta evanescente.

\end{proof}

\begin{remark}

El resultado es específico de la subsecuencia convergente seleccionada.

El costo normalizado:

\[
\frac{G_B(K)}{K}
\]

puede disminuir para muchas secuencias de tamaños de bloque
crecientes.

Sin embargo, el crecimiento arbitrario de $K$ no garantiza que la
brecha de redondeo entero absoluta disminuya.

RICAN proporciona una regla de selección determinista mediante
fracciones continuas que identifica longitudes de bloque con error de
redondeo absoluto excepcionalmente pequeño.

\end{remark}

% =================================================
% V. ALGORITMO DE SELECCIÓN DE TAMAÑO DE BLOQUE
% =================================================

\section{Algoritmo de Selección de Tamaño de Bloque}

Los resultados teóricos proporcionan un procedimiento constructivo
para seleccionar tamaños de bloque RICAN.

Dada una cardinalidad de alfabeto $B$ y un umbral de brecha de
representación deseado $\epsilon$, el algoritmo calcula los
convergentes de fracciones continuas de:

\[
\log_2(B).
\]

Cada denominador se evalúa como un tamaño de bloque candidato.

Para cada candidato:

\[
K=q_i,
\]

se calcula la longitud binaria requerida:

\[
n(K)=\lceil K\log_2(B)\rceil .
\]

La brecha de representación es:

\[
G_B(K)
=
n(K)-K\log_2(B).
\]

Para el alfabeto decimal:

\[
B=10,
\]

la implementación de referencia utiliza:

\[
K=643 .
\]

Este tamaño de bloque proporciona un equilibrio práctico entre
complejidad aritmética y ajuste de la representación.

Valores derivados de convergentes mayores, tales como:

\[
K=783032,
\]

demuestran el comportamiento asintótico de la construcción pero
requieren operaciones aritméticas con enteros más grandes.

Todas las evaluaciones numéricas de convergentes deben utilizar
aritmética de precisión arbitraria para evitar errores de redondeo
de punto flotante.

\begin{algorithm}[!t]
\caption{Selección de Bloque RICAN Usando Fracciones Continuas}
\label{alg:block}
\begin{algorithmic}[1]
\REQUIRE Cardinalidad de alfabeto $B\geq2$, brecha objetivo $\epsilon$
\ENSURE Tamaño de bloque $K$, longitud binaria $n(K)$
\STATE Calcular:
\[
x\leftarrow\log_2(B)
\]
\STATE Generar expansión en fracciones continuas de $x$
\STATE Calcular convergentes:
\[
p_i/q_i
\]
\FOR{cada convergente}
\STATE
\[
K\leftarrow q_i
\]
\STATE
\[
n(K)\leftarrow\lceil Kx\rceil
\]
\STATE
\[
G\leftarrow n(K)-Kx
\]
\IF{$G<\epsilon$}
\RETURN $(K,n(K),G)$
\ENDIF
\ENDFOR
\end{algorithmic}
\end{algorithm}

% =================================================
% VI. MARCO DE REPRESENTACIÓN
% =================================================

\section{Marco de Representación}

\subsection{Modelo General de Codificación}

El proceso de codificación RICAN construye un mapeo inyectivo desde
bloques de alfabeto finito a palabras binarias.

Dada una secuencia de entrada:

\[
S=s_0s_1\ldots s_{L-1}
\]

sobre un alfabeto:

\[
\Sigma_B
\]

con cardinalidad $B$, y un tamaño de bloque seleccionado $K$, la
secuencia se particiona en bloques.

Para un bloque completo:

\[
B_i=(s_0,s_1,\ldots,s_{K-1}),
\]

el valor entero correspondiente es:

\[
X(B_i)
=
\sum_{j=0}^{K-1}
s_j B^{K-1-j}.
\]

Los valores posibles satisfacen:

\[
0\leq X(B_i)<B^K .
\]

Dado que:

\[
B^K\leq2^{n(K)},
\]

donde:

\[
n(K)=\lceil K\log_2(B)\rceil ,
\]

cada bloque puede representarse de manera única usando exactamente
$n(K)$ bits.

El bloque incompleto final, cuando está presente, se codifica usando
su longitud real:

\[
K'<K .
\]

El decodificador reconstruye la secuencia original aplicando el mapeo
inverso de entero a símbolo usando división repetida por $B$.

\subsection{Mapeo de Representación Inyectiva}

Para un alfabeto y tamaño de bloque fijos, el codificador se define
como:

\[
E_B:
\Sigma_B^L
\rightarrow
\{0,1\}^{N(L)} .
\]

El mapeo es inyectivo porque cada bloque de entrada válido se asigna
a un valor entero único.

El decodificador es el mapeo inverso sobre la imagen del codificador:

\[
D_B:
\operatorname{Im}(E_B)
\rightarrow
\Sigma_B^L .
\]

Por lo tanto:

\[
D_B(E_B(S))=S
\]

para toda secuencia de entrada válida $S$.

El espacio de representación binaria puede contener estados no
utilizados porque:

\[
2^{n(K)}>B^K
\]

para alfabetos no potencia de dos.

Estos estados no utilizados no afectan la reversibilidad porque el
decodificador opera únicamente sobre representaciones codificadas
válidas.

\subsection{Alfabetos de Potencia de Dos}

Para:

\[
B=2^m ,
\]

la representación es exacta:

\[
n(K)
=
\lceil Km\rceil
=
Km .
\]

Por lo tanto:

\[
G_B(K)=0 .
\]

Ejemplos incluyen:

\begin{itemize}

\item Alfabeto binario ($B=2$): un bit por símbolo.

\item Alfabeto de ADN ($B=4$): dos bits por símbolo.

\item Alfabeto hexadecimal ($B=16$): cuatro bits por símbolo.

\item Alfabeto de bytes ($B=256$): ocho bits por símbolo.

\end{itemize}

\subsection{El Formato de Representación Autónomo RJ03}

La implementación de referencia utiliza el formato de representación
binaria RJ03.

RJ03 está diseñado como un contenedor de representación reversible
autónomo.

A diferencia de los formatos codificados entrópicamente, RJ03 no
requiere:

\begin{itemize}

\item modelos de probabilidad,

\item estados adaptativos,

\item diccionarios externos,

\item datos de entrenamiento.

\end{itemize}

Todos los parámetros requeridos para la reconstrucción se almacenan
dentro del archivo de representación.

El formato contiene:

\begin{itemize}

\item definición del alfabeto,

\item cardinalidad del alfabeto,

\item tamaño de bloque,

\item longitud original de la secuencia,

\item flujo de bloques codificados.

\end{itemize}

Los campos binarios utilizan representación little-endian sin signo.

Los bits no utilizados en el byte final son bits de relleno y se
ignoran durante la decodificación.

\begin{table}[!t]
\caption{Formato de Representación Autónomo RJ03}
\label{tab:rj03}
\centering
\begin{tabular}{lll}
\toprule
Desplazamiento & Tamaño & Campo\\
\midrule
0 & 4 & Identificador mágico ``RJ03''\\
4 & 4 & Versión del formato\\
8 & 4 & Cardinalidad del alfabeto $B$\\
12 & 8 & Tamaño de bloque $K$\\
20 & 8 & Bits por bloque $n(K)$\\
28 & 4 & Longitud del alfabeto\\
32 & variable & Símbolos del alfabeto\\
-- & 8 & Longitud de secuencia original $L$\\
-- & 8 & Conteo de bloques completos\\
-- & 8 & Tamaño de bloque residual\\
-- & variable & Flujo de bits codificado\\
\bottomrule
\end{tabular}
\end{table}

% =================================================
% VII. EVALUACIÓN EXPERIMENTAL
% =================================================

\section{Evaluación Experimental}

Los experimentos evalúan las propiedades matemáticas de representación
de RICAN.

El objetivo no es la comparación de compresión.

Los experimentos verifican:

\begin{enumerate}

\item reconstrucción exacta,

\item concordancia entre tamaños de representación teóricos y
prácticos,

\item independencia de la longitud de representación de las
estadísticas de la fuente,

\item aplicabilidad a múltiples alfabetos finitos.

\end{enumerate}

\subsection{Verificación de Representación Decimal}

Una secuencia decimal que contiene:

\[
L=1\,000\,508
\]

dígitos fue codificada con:

\[
B=10
\]

y:

\[
K=643 .
\]

El tamaño de bloque teórico es:

\[
n(K)=2136
\]

bits.

El tamaño medido del contenedor se muestra en la
Tabla~\ref{tab:decimal-validation}.

\begin{table}[!t]
\caption{Verificación de representación decimal}
\label{tab:decimal-validation}
\centering
\begin{tabular}{lc}
\toprule
Parámetro & Valor\\
\midrule
Símbolos de entrada & 1,000,508\\
Tamaño del alfabeto & 10\\
Tamaño de bloque & 643\\
Bits por bloque & 2136\\
Tamaño codificado & 415,474 bytes\\
Verificación & PASO\\
\bottomrule
\end{tabular}
\end{table}

El tamaño medido incluye el costo del contenedor RJ03 requerido para
la reconstrucción autónoma.

\subsection{Independencia de la Longitud de Representación}

Debido a que RICAN es un marco de representación de tasa fija, el
tamaño de salida depende de:

\[
B
\]

y:

\[
L .
\]

No depende de la distribución de probabilidad de los símbolos.

Secuencias con diferentes propiedades estadísticas pero longitud
idéntica y cardinalidad de alfabeto idéntica producen tamaños de
representación idénticos.

\begin{table}[!t]
\caption{Independencia de la Longitud de Representación}
\label{tab:independence}
\centering
\begin{tabular}{lc}
\toprule
Tipo de Fuente & Tamaño Codificado\\
\midrule
Dígitos de constantes matemáticas & Igual\\
Secuencia decimal aleatoria & Igual\\
Secuencia decimal estructurada & Igual\\
\bottomrule
\end{tabular}
\end{table}

\subsection{Generalización de Alfabetos}

El mismo marco se aplica a múltiples alfabetos finitos.

\begin{table}[!t]
\caption{Verificación de Alfabetos Finitos}
\label{tab:alphabet}
\centering
\begin{tabular}{ccc}
\toprule
Alfabeto & $B$ & Resultado\\
\midrule
Binario & 2 & PASO\\
ADN & 4 & PASO\\
Decimal & 10 & PASO\\
Hexadecimal &16& PASO\\
ASCII &128& PASO\\
\bottomrule
\end{tabular}
\end{table}

% =================================================
% VIII. DISCUSIÓN
% =================================================

\section{Discusión}

\subsection{Contribución Principal}

La contribución principal de RICAN es la construcción de esquemas de
representación de tasa fija cuyo costo de redondeo entero puede
reducirse mediante selección de bloques por fracciones continuas.

La existencia de representaciones de longitud fija se sigue
directamente de la cardinalidad de espacio finito.

RICAN contribuye con un método determinista para seleccionar tamaños
de bloque con brechas de representación excepcionalmente pequeñas.

\subsection{Relación con la Compresión}

RICAN no reduce la información de la fuente.

No explota:

\begin{itemize}

\item frecuencias de símbolos,

\item patrones repetidos,

\item dependencias estadísticas.

\end{itemize}

Por lo tanto, RICAN no debe considerarse un reemplazo para métodos de
codificación entrópica.

Su objetivo es la optimización del costo determinista de
representación binaria.

\subsection{Limitaciones}

El marco actual tiene varias limitaciones.

Los tamaños de bloque convergentes muy grandes requieren aritmética
de enteros de precisión arbitraria.

El formato de representación introduce un costo de metadatos
constante.

Además, debido a que el método no utiliza modelos de probabilidad, no
proporciona ganancias de compresión para fuentes redundantes.

% =================================================
% IX. DISPONIBILIDAD DEL CÓDIGO
% =================================================

\section{Disponibilidad del Código y Reproductibilidad}

La implementación de referencia, la especificación RJ03 y los scripts
de evaluación están disponibles públicamente en:

\begin{center}
\url{https://github.com/SequenceOfManyZerosAndOnes/RICAN}
\end{center}

El repositorio contiene:

\begin{itemize}

\item implementación del codificador y decodificador,

\item documentación del formato RJ03,

\item ejemplos reproducibles,

\item documentación matemática.

\end{itemize}

Los experimentos se realizaron utilizando:

\begin{verbatim}
Python 3.13
RICAN RJ03 referencia de implementación
\end{verbatim}

Se planea un lanzamiento de archivo permanente a través de Zenodo.

Versión del software:

\[
\text{RICAN RJ03 v1.0.0 -- Julio 2026}
\]

% =================================================
% X. CONCLUSIÓN
% =================================================

\section{Conclusión}

Este artículo presentó RICAN, un marco de representación binaria de
tasa fija para alfabetos finitos basado en selección de convergentes
de fracciones continuas.

El marco se aplica uniformemente a:

\[
\mathcal{A}
=
\{\Sigma_B:B\geq2\}.
\]

Para alfabetos de potencia de dos:

\[
B=2^m,
\]

la representación es exacta.

Para alfabetos no potencia de dos, los convergentes de fracciones
continuas de $\log_2(B)$ proporcionan longitudes de bloque cuya
brecha de representación absoluta se aproxima a cero a lo largo de la
subsecuencia seleccionada.

Para la representación decimal:

\[
B=10,
\]

la implementación práctica utiliza:

\[
K=643,
\]

mientras que valores derivados de convergentes mayores demuestran el
comportamiento asintótico de la construcción.

RICAN proporciona un marco matemático para estudiar el límite entre
la representación de alfabetos finitos y el almacenamiento binario.

Las direcciones de investigación futura incluyen:

\begin{itemize}

\item representación de ADN a gran escala,

\item aceleración por hardware,

\item implementaciones paralelas,

\item optimización de bloques convergentes muy grandes.

\end{itemize}

% =================================================
% REFERENCIAS
% =================================================

\bibliographystyle{IEEEtran}
\begin{thebibliography}{00}

\bibitem{shannon1948}
C. E. Shannon,
``Una Teoría Matemática de la Comunicación,''
\emph{Bell System Technical Journal},
vol. 27,
pp. 379--423, 623--656,
1948.

\bibitem{huffman1952}
D. A. Huffman,
``Un Método para la Construcción de Códigos de Redundancia Mínima,''
\emph{Proceedings of the IRE},
vol. 40,
no. 9,
pp. 1098--1101,
1952.

\bibitem{ziv1977}
J. Ziv y A. Lempel,
``Un Algoritmo Universal para Compresión de Datos Secuenciales,''
\emph{IEEE Transactions on Information Theory},
vol. 23,
no. 3,
pp. 337--343,
1977.

\bibitem{rissanen1979}
J. Rissanen y G. G. Langdon,
``Codificación Aritmética,''
\emph{IBM Journal of Research and Development},
vol. 23,
no. 2,
pp. 149--162,
1979.

\bibitem{duda2013}
J. Duda,
``Sistemas Numéricos Asimétricos: Codificación Entrópica que Combina
Velocidad de Huffman con Compresión de Codificación Aritmética,''
arXiv:1311.2540,
2013.

\bibitem{cover1973}
T. M. Cover,
``Codificación Enumerativa de Fuente,''
\emph{IEEE Transactions on Information Theory},
vol. 19,
no. 1,
pp. 73--77,
1973.

\bibitem{immink2004}
K. A. S. Immink,
``Una Encuesta de Técnicas de Codificación Enumerativa,''
\emph{IEEE Transactions on Information Theory},
vol. 50,
no. 10,
pp. 2276--2285,
2004.

\bibitem{hardy2008}
G. H. Hardy y E. M. Wright,
\emph{Una Introducción a la Teoría de Números},
6ta ed.,
Oxford University Press,
2008.

\bibitem{khinchin1997}
A. Ya. Khinchin,
\emph{Fracciones Continuas},
Dover Publications,
1997.

\bibitem{cassels1957}
J. W. S. Cassels,
\emph{Una Introducción a la Aproximación Diofántica},
Cambridge University Press,
1957.

\bibitem{rockett1992}
A. M. Rockett y P. Sz\"usz,
\emph{Fracciones Continuas},
World Scientific,
1992.

\bibitem{brent2010}
R. P. Brent y P. Zimmermann,
\emph{Aritmética Computacional Moderna},
Cambridge University Press,
2010.

\end{thebibliography}

\end{document}
